\scanpage{103}
% Archive foliation: 104, inverted at upper right. Reused unrelated typescript; printed page 47.
\begin{center}-- 47 --\hfill n\textsuperscript{o} 380\end{center}
\begin{center}\S\ 4. \uuline{Algèbres de fonctions}.\end{center}
1. \uline{Algèbres de fonctions continues bornées}.

Etant donné un espace topologique $\Omega$, on note $\mathcal B(\Omega)$ l'algèbre sur $\mathbb C$ des fonctions \inserted{continues} bornées sur $\Omega$ à valeurs dans $\mathbb C$. C'est une algèbre commutative semi-simple. L'application de $\mathcal B(\Omega)$ dans $\mathcal B(\Omega)$ qui transforme chaque fonction $f$ en la fonction conjuguée $\bar f$ est une involution de $\mathcal B(\Omega)$. On définit une norme sur $\mathcal B(\Omega)$ en posant $\lVert f\rVert=\underset{t\in\Omega}{\operatorname{Sup}}|f(t)|$ pour tout $f\in\mathcal B(\Omega)$. Cette norme est compatible avec la structure d'algèbre. Munie de cette norme, $\mathcal B(\Omega)$ est une algèbre normée complète. L'involution $f\longrightarrow\bar f$ est une isométrie. Pour tout élément $f$ hermitien, c'est-à-dire pour toute fonction bornée continue réelle, $1+f^2$ est inversible dans $\mathcal B(\Omega)$. Par suite la transformation de Gelfand de $\mathcal B(\Omega)$ vérifie la condition $\widehat{\bar f}=\overline{\hat f}$. (Prop.10). Il en résulte que l'image de la transformée de Gelfand de $\mathcal B(\Omega)$ est une sous-algèbre dense dans $\mathcal C(\operatorname{Sp}(\mathcal B(\Omega)))$ (Prop.10, Cor.\quad). Mais d'autre part, on a $\lVert f^2\rVert=\lVert f\rVert^2$ pour tout $f\in\mathcal B(\Omega)$ et par suite la transformation de Gelfand est une isométrie. On voit donc que la transformation de Gelfand est un isomorphisme isométrique de $\mathcal B(\Omega)$ sur $\mathcal C(\operatorname{Sp}(\mathcal B(\Omega)))$.

Si $\varphi$ est une application continue de l'espace $\Omega$ dans un espace $\Omega'$, on définit un homomorphisme $\varphi^{\mathcal B}$ de $\mathcal B(\Omega')$ dans $\mathcal B(\Omega)$ en posant $\varphi^{\mathcal B}(f)=f\circ\varphi$ pour tout $f\in\mathcal B(\Omega')$. Cet homomorphisme est continu car $\lVert\varphi^{\mathcal B}(f)\rVert\leq\lVert f\rVert$ pour tout $f\in\mathcal B(\Omega')$.

Pour tout $t\in\Omega$, l'application $\xi_t$ de $\mathcal B(\Omega)$ dans $\mathbb C$ définie par $\xi_t(f)=f(t)$ quel que soit $f\in\mathcal B(\Omega)$ est un caractère de $\mathcal B(\Omega)$. L'application $t\longrightarrow\xi_t$ est une application continue qui sera appelée l'application canonique de $\Omega$ dans le spectre de $\mathcal B(\Omega)$.

Pour que cette application soit injective, il faut et il suffit que $\mathcal B(\Omega)$
